\providecommand{\toprule}{\hline\hline}
\providecommand{\midrule}{\hline}
\providecommand{\bottomrule}{\hline\hline}
\providecommand{\addlinespace}[1][]{}
{\small
\renewcommand{\arraystretch}{1.12}
\begin{tabular}{llrrrr}
\toprule
Coorte & Período & Voltou 6m: s/ seguro (<6m) & Voltou 6m: seguro 5p (>=24m) & $\Delta\ln w$: Meses parado & $\Delta\ln w$: Meses $\times$ Portab. \\
\midrule
2014 & Véspera da recessão & -0,0212 (0,0155) & -0,0065 (0,0126) & -0,0034*** (0,0009) & -0,0070*** (0,0015) \\
2015 & Recessão & -0,0253 (0,0164) & -0,0093 (0,0108) & -0,0032*** (0,0007) & -0,0055*** (0,0011) \\
2016 & Recessão & -0,0291 (0,0182) & -0,0073 (0,0119) & -0,0026*** (0,0006) & -0,0032*** (0,0007) \\
2019 & Véspera da pandemia & -0,0144 (0,0163) & +0,0045 (0,0104) & -0,0085*** (0,0006) & -0,0020** (0,0008) \\
2020 & Pandemia & +0,0137 (0,0121) & +0,0095 (0,0121) & -0,0001 (0,0005) & +0,0001 (0,0006) \\
2021 & Pandemia & -0,0135 (0,0175) & +0,0010 (0,0134) & -0,0022*** (0,0004) & +0,0012** (0,0005) \\
\bottomrule
\end{tabular}
\vspace{2pt}
\par{\footnotesize \textit{Nota:} As duas primeiras colunas comparam o efeito da portabilidade sobre o retorno em até 6 meses entre demitidos sem elegibilidade ao seguro-desemprego (tenure < 6 meses) e com elegibilidade máxima (tenure >= 24 meses, 5 parcelas). As duas últimas colunas estimam o gradiente entre variação salarial e duração do desemprego interagido com a portabilidade.}
}
